\section{Related Work}
\label{sec:related}

Several research threads intersect in this study: compound-fault
decoupling, zero-shot recognition of unseen fault combinations, the
reliability of evaluation protocols in data-driven condition monitoring,
generalization across operating conditions, multimodal motor diagnosis, and
the benchmark datasets that underpin them. We review each and then position
the present work (Table~\ref{tab:positioning}).

Early data-driven treatments of compound faults assigned each observed
combination its own class, which cannot express combinations absent from
training. Decoupling formulations address this by predicting the elementary
mechanisms instead. Huang \emph{et al.} proposed a deep decoupling
convolutional network whose multi-label classifier recovers the
constituents of gearbox compound faults from raw vibration signals
\cite{huang2019ddcnn}, and a one-dimensional CNN variant with a multi-label
output layer \cite{huang2019multilabel}. Subsequent work refined the
decoupling head: capsule networks with dynamic routing and threshold-based
label activation \cite{ltssbow2024}, attentional residual networks combined
with active learning for measured (rather than injected) bearing compound
faults \cite{decoupling_attn2020}, and mixture-of-experts architectures
with task-adapted attention for decoupling compounds unseen as combinations
\cite{hemtan2024}. Two properties are near-universal in this line: the
constituents are \emph{mechanical} (bearing races, balls, gear teeth), so
every constituent is observable in the same vibration channel; and
evaluation is sample-level on steady-speed laboratory recordings, usually
with segments of the same recording on both sides of the train/test split.

A stricter setting withholds \emph{all} compound data and requires
composing single-fault knowledge at test time. Xu \emph{et al.} introduced
zero-shot compound-fault diagnosis with manually defined semantic
attributes learned from single-fault data \cite{eswa2022}, later replacing
hand-crafted attributes with autoencoder-learned fault semantics
\cite{eswa2023}. Generative formulations synthesize pseudo-samples of
unseen compounds from single-fault semantics \cite{zsl2021,kbs2025}; graph
embeddings propagate relational structure among base faults
\cite{neuro2025}; and physics-informed semantic spaces anchor the attribute
vectors in characteristic-frequency knowledge \cite{song2025gzsl}. Reported
accuracies on unseen compounds cluster between roughly 75\% and 87\%
\cite{eswa2022,eswa2023,song2025gzsl,appac2025}.

These results, however, concern compounds whose constituents are
\emph{homogeneous}: bearing--bearing or bearing--gear combinations in which
both mechanisms modulate the same vibration carrier. The compounds in the
MCC5-THU motor dataset instead pair a mechanical bearing defect with an
electrical or electromagnetic fault (winding short circuit, broken rotor
bars, eccentricity), whose primary evidence lies in the stator currents
rather than in vibration. To our knowledge, this \emph{cross-modal}
compound regime has not been evaluated in the zero-shot compound-fault
literature, and no published constituent-level results exist for this
dataset.

A parallel literature shows that evaluation-protocol choices, more than
model choices, drive many reported near-ceiling accuracies. Wheat
\emph{et al.} quantified the impact of intra-bearing data leakage in
vibration diagnosis and advocated part-to-part splits \cite{wheat2024};
recent work formalizes the distinction between segmentation-level and
specimen-level leakage and proposes leakage-safe protocols
\cite{gama2025,phme2025}; and open benchmark studies made similar
observations for deep architectures across public datasets
\cite{zhao2020benchmark}. This thread has concentrated on single-fault
classification. Its intersection with compound-fault evaluation is, to our
knowledge, unexplored---yet it is precisely where leakage is most subtle:
in the dataset used here, severity variants of one physical combination can
leak across a composition boundary, and compound recordings are acquired in
different sessions from their single-fault references, which confounds
naive ``compound versus single'' contrasts. Our protocol design (run-level,
severity-aware splits; session-matched sensitivity controls;
constituent-prevalence floors for every reported metric) extends
leakage-aware evaluation to the compositional setting.

Diagnosis under variable speed and load has produced a large
domain-adaptation and domain-generalization literature, surveyed in
\cite{han2025}. On the MCC5-THU family specifically, the dataset authors
proposed evidential ensemble learning for real-time multimode diagnosis
\cite{liu2024tii} and recently questioned how operating conditions should
enter learning-based diagnosis at all \cite{rethinking2025}; other groups
applied condition-robust architectures to the companion gearbox release
\cite{wang2024,phmap2025}. The dominant protocol in this literature holds
one condition out of many \cite{han2025,domainshift2025}. Our reference
experiments on this dataset (Appendix~\ref{app:condition}) show that
protocol to be nearly saturated---training on eleven of twelve conditions
costs only ${\sim}3$ accuracy points, because the remaining conditions
bracket the held-out one, whereas the operationally realistic reverse
direction (training on a single condition) costs over 60---which motivates
the crossed composition--context--condition design used here rather than
leave-one-condition-out alone.

That mechanical faults imprint on stator currents and electrical faults on
vibration has long motivated multimodal motor diagnosis
\cite{mcsa_review2017}. Deep multimodal fusion typically concatenates
modalities into one representation or fuses learned features, e.g.\
multilevel information fusion \cite{multilevel2019}, dynamic-routing
multimodal networks \cite{drmnn2020}, adaptive CNN fusion of denoised
vibration and current \cite{plos2025}, and comparative studies of the two
modalities under deep and classical learners \cite{ayankoso2026}. These
works report that fusion improves \emph{single-fault} accuracy. The
compound setting reverses the picture: the reference experiments of
Appendix~\ref{app:condition} show that a shared feature space lets the
dominant (mechanical) signature suppress the electrical one---a winding
short detectable from currents alone on 39.3\% of compound windows falls to
0.2\% once vibration features are concatenated. The mechanism-specific feature families used here are a
structural answer: each constituent detector sees only the physically
appropriate modality, so cross-modal suppression cannot occur inside a
detector.

Public rotating-machinery datasets underpin this literature: CWRU
\cite{smith2015cwru} and Paderborn \cite{lessmeier2016} for bearings, and
more recent multi-fault motor and drivetrain releases such as a PMSM
stator-fault dataset \cite{jung2023} and an industrial-scale motor--pump
dataset \cite{bruinsma2024}. The MCC5-THU family---a gearbox release
\cite{chen2024gearbox} and the motor release used here \cite{chen2026}---is
distinctive in combining single \emph{and} compound faults, twelve
speed/load profiles with transitional segments, and synchronized vibration,
three-phase current, torque, and key-phase channels. The motor release is
distributed without reference baselines; the present study provides, to our
knowledge, its first systematic constituent-level evaluation, together with
documentation of release-level inconsistencies that materially affect
protocol design.

Table~\ref{tab:positioning} contrasts this study with representative prior
work along the axes developed above. The distinguishing combination is:
compound faults whose constituents live in \emph{different} modalities;
composition, compound-context, and operating-condition generalization
varied \emph{jointly} in one crossed protocol; run-level, severity-aware,
leakage-safe splits; metrics reported against constituent-prevalence floors
so that prior-driven scores cannot masquerade as diagnostic skill; and
run-clustered bootstrap inference with multiplicity control.

\begin{table*}[!t]
\caption{Positioning of this study against representative prior work.
$\checkmark$~=~explicitly included; --~=~absent, not applicable, or not
reported. ``Cross-modal'' means the compound's constituents are observable
in different sensing modalities (mechanical $+$ electrical). ``Prior-aware
floors'' means metrics are reported against label-prevalence (trivial
constant-predictor) references.}
\label{tab:positioning}
\centering
\footnotesize
\setlength{\tabcolsep}{4.5pt}
\begin{tabular}{l l l c c c c c c}
\toprule
Study & Task focus & Signals & \makecell{Compound\\faults} &
\makecell{Cross-\\modal} & \makecell{Unseen\\combin.} &
\makecell{Condition\\shift} & \makecell{Leakage-safe\\run-level splits} &
\makecell{Prior-aware\\floors} \\
\midrule
Huang \emph{et al.} \cite{huang2019ddcnn} & multi-label decoupling & vib. & $\checkmark$ & -- & -- & -- & -- & -- \\
LTSS-BoW-CapsNet \cite{ltssbow2024} & capsule decoupling & vib. & $\checkmark$ & -- & -- & -- & -- & -- \\
HeMTAN \cite{hemtan2024} & unseen-compound decoupling & vib. & $\checkmark$ & -- & $\checkmark$ & -- & -- & -- \\
Xu \emph{et al.} \cite{eswa2022} & zero-shot compounds & vib. & $\checkmark$ & -- & $\checkmark$ & -- & -- & -- \\
Xu \emph{et al.} \cite{eswa2023} & learned fault semantics & vib. & $\checkmark$ & -- & $\checkmark$ & -- & -- & -- \\
Song \emph{et al.} \cite{song2025gzsl} & GZSL, physics semantics & vib. & $\checkmark$ & -- & $\checkmark$ & -- & -- & -- \\
Wheat \emph{et al.} \cite{wheat2024} & leakage in evaluation & vib. & -- & -- & -- & -- & $\checkmark$ & -- \\
Han \emph{et al.} \cite{han2025} & multi-condition survey & multi & -- & -- & -- & $\checkmark$ & -- & -- \\
Liu \emph{et al.} \cite{liu2024tii} & multimode diagnosis & multi & -- & -- & -- & $\checkmark$ & -- & -- \\
DRMNN \cite{drmnn2020} & multimodal fusion & vib.+cur. & -- & -- & -- & -- & -- & -- \\
\midrule
\textbf{This work} & \makecell[l]{constituent-level\\generalization} &
\makecell[l]{vib.+cur.\\(+key-phase)} & $\checkmark$ & $\checkmark$ &
$\checkmark$ & $\checkmark$ & $\checkmark$ & $\checkmark$ \\
\bottomrule
\end{tabular}
\end{table*}
